% ===================================================================
%  LENTELĖ Nr. 8 — Derliaus nuėmimas.
%  Traduction 2026 d'apres texto/io, controlee sur texto/fr, texto/ru
%  et texto/cs. Consigne : voir texto/lt/10-lentele-01.tex.
%
%  L'ORAGE DE L'ALINEA 12 DONNE A CETTE COLONNE UN MOT QU'AUCUNE DES
%  TREIZE PRECEDENTES NE POUVAIT DONNER. Le tonnerre se dit
%  « perkūnas », l'orage « perkūnija » : c'est le nom du DIEU du
%  tonnerre de la religion d'avant, reste mot courant et sans
%  majuscule dans la langue de tous les jours. Le tcheque dit
%  « hrom », le galicien « trono », l'irlandais « toirneach » — des
%  mots de bruit. Le lituanien dit le nom d'un dieu, parce que la
%  Lituanie a ete christianisee la derniere de l'Europe et que le mot
%  a eu le temps de devenir ordinaire.
%
%  ET LES DEUX AUTRES MOTS DE L'ORAGE DISENT ENCORE L'ACTION, comme
%  aux tableaux 6 et 7 : « žaibas » l'eclair et « griaustinis » le
%  grondement, de « griauti », abattre. Le meme alinea porte donc les
%  deux couches de la langue cote a cote — le tres vieux nom propre
%  devenu nom commun, et le nom commun fabrique sur un verbe.
%
%  LE BLEUET S'APPELLE ICI LA FLEUR DU SEIGLE. « rugiagėlė », de
%  « rugiai » le seigle et « gėlė » la fleur : la plante est nommee
%  par le champ ou elle pousse, comme le cheval l'etait par l'outil
%  qu'il tire. Le tableau 7 avait montre l'outil nommant la bete ; le
%  tableau 8 montre le champ nommant la fleur. C'est la meme habitude
%  de langue, prise deux fois sur deux pages voisines.
%
%  ET « aciani » SONT DES BLEUETS, NON DU BLEU. Le releve des
%  couleurs de l'ido est trompeur pour qui va vite : « aciano » est
%  la fleur, le nom de la couleur n'en est qu'un derive. La colonne
%  galicienne s'y etait deja reprise ; on le note ici pour que la
%  prochaine ne s'y reprenne pas.
%
%  L'ECART DECLARE DU BLOC c2-03-1, LE DEUXIEME DES TROIS. Le
%  fac-simile ido y ecrit « gobleto 13) », sans la parenthese
%  ouvrante. UNE TRADUCTION N'EST PAS UNE TRANSCRIPTION : la coquille
%  appartient au compositeur de l'ido, non au texte lituanien, et les
%  treize colonnes precedentes ont toutes retabli la parenthese. On
%  fait comme elles ; renvoji.py passe l'ecart sans rien dire.
% ===================================================================

%%K t08-noto-1 noto
\VUnotes{143.2mm}{%
\textit{(*) Kadangi tas žodis nėra tikslus: ar tai linų nuėmimas, vynuogių
nuėmimas, obuolių nuėmimas? Gal būtų gera vartoti „} \VUgras{mesono}
\textit{“, pasiūlytą Mondo XIV, p.~242. Bet iki šiol ta nauja šaknis nėra Akademijos
priimta ir laukia svarstymų pagal įprastas idistų taisykles.}%
}

%%K t08-apar-1 apar
\VUsaut{5.0mm}
\VUtitre{15.0pt}{60}{LENTELĖ Nr. 8}
\VUsaut{3.0mm}
\VUcentre{13.2pt}{90}{\textit{Pirmoji scena.}}
\VUsaut{3.0mm}
\VUpk{10.2pt}{40}{Derliaus nuėmimas (*).}
\VUsaut{4.0mm}
\VUpk{10.2pt}{40}{Kaimo vaizdai.}
\VUsaut{3.0mm}

%%K t08-c1-01-1 p
1. --- Su \VUgras{vasara} dar kartą pasikeitė kaimo vaizdas. Vagai patikėta sėkla
sudygo; iš žemės išlindo žalias daigelis ir išleido varpą. Grūdas susiformavo, paskui
subrendo, ir nuo šiol belieka tik nuimti derlių.

%%K t08-c1-02-1 p
2. --- \VUgras{Laukų gėlės} pražydo. \VUgras{Rugiagėlės} \textsuperscript{(1)},
\VUgras{aguonos} \textsuperscript{(2)} ir \VUgras{rusmenės} \textsuperscript{(3)}
įmaišo į vienodą kviečių laukų atspalvį linksmą savo šviežių \VUgras{vainikėlių}
kontrastą.

%%K t08-c1-03-1 p
3. --- Vaismedžiai apsipylė lapais; vaisiai subrendo. Mažoji \VUgras{vyšnia}
\textsuperscript{(4)} nusėta gražiomis \VUgras{vyšniomis} \textsuperscript{(5)},
kurias vaikai su didžiausiu malonumu skina ir valgo.

%%K t08-c1-04-1 p
4. --- Nuo šiol gyvuliai gali visą dieną praleisti gryname ore.

%%K t08-c1-04-2 p
\VUgras{Ėriukai} \textsuperscript{(6)} užaugo ir virto gražiomis \VUgras{avimis}
\textsuperscript{(7)} ir gražiais \VUgras{avinais} \textsuperscript{(8)} su tankiu
vilnos kailiu. Jie ramiai ganosi \VUgras{žolę} \VUgras{ganykloje}
\textsuperscript{(9)}, marginamoje \VUgras{ramunėlėmis.}

%%K t08-c1-04-3 p
Netrukus tie gyvuliai bus kerpami, kad būtų gauta jų \VUgras{vilna,} iš kurios daroma
mūsų drabužių gelumbė.

%%K t08-tit-2 sub
\VUsaut{5.0mm}
\VUpk{10.2pt}{40}{Piemuo.}
\VUsaut{3.4mm}

%%K t08-c1-05-1 p
5. --- \VUgras{Piemuo} \textsuperscript{(10)} tarsi nelabai į juos žiūri. Jam rūpi
tik audra, kuri eina ties horizontu, o \VUgras{avių sargas} \textsuperscript{(13)},
kuris kelia \VUgras{gertuvę} \textsuperscript{(14)} prie lūpų, atrodo dar
nerūpestingesnis.

%%K t08-c1-06-1 p
6. --- Karštis slegiantis; mat ir \VUgras{šuo} \textsuperscript{(15)} ateina
numalšinti troškulio; jis laka gaivų \VUgras{upelio} \textsuperscript{(16)} vandenį ir
išgąsdina vargšę \VUgras{varlytę} \textsuperscript{(17)}, kuri buvo pritūpusi kranto
žolėse. Ši šoka į skaidrų vandenį, kur žydi \VUgras{vandens lelijos}
\textsuperscript{(18)}.

%%K t08-c1-07-1 p
7. --- \VUgras{Kielė} \textsuperscript{(19)} maudosi prie \VUgras{šaltinėlio}
\textsuperscript{(20)}, kuris trykšta tarp dviejų uolų, ir slepiasi po
\VUgras{dygliuotais krūmais} \textsuperscript{(21)} ir \VUgras{gudobelių krūmais}
\textsuperscript{(22)}.

%%K t08-tit-3 sub
\VUsaut{5.0mm}
\VUpk{10.2pt}{40}{Derliaus nuėmimas.}
\VUsaut{3.4mm}

%%K t08-c1-08-1 p
8. --- Kaimo vaikams kviečių pjūtis yra labai linksmas metas. Jie džiaugsmingai
\VUgras{verčiasi kūliais} \textsuperscript{(24)} ant \VUgras{šiaudų kūgių}
\textsuperscript{(23)}, jie padeda \VUgras{pjovėjams} \textsuperscript{(26)}, ir kol
šie pjauna \VUgras{kviečių lauką} \textsuperscript{(27)} savo \VUgras{dalgiais}
\textsuperscript{(25)}, gerai išgaląstais \VUgras{galąstuvu} \textsuperscript{(30)},
jie renka kviečius ir riša juos į \VUgras{pėdus} \textsuperscript{(31)} arba
\VUgras{glėbelius} \textsuperscript{(32)}.

%%K t08-c1-09-1 p
9. --- Kartais didžiausiems iš jų leidžiama \VUgras{pjautuvu} \textsuperscript{(12)}
pjauti \VUgras{auksinius stiebus} \textsuperscript{(28)}, kurie neša sunkias
\VUgras{varpas} \textsuperscript{(29)}, pilnas grūdų; o mažieji, saugodami
\VUgras{bandą} \textsuperscript{(33)}, kuri ganosi \VUgras{ganykloje}
\textsuperscript{(34)} didžiosios \VUgras{liepos} \textsuperscript{(35)} pavėsyje,
gaudo savo \VUgras{tinkleliu} \textsuperscript{(36)} gražius \VUgras{drugelius.}

%%K t08-c1-09-2 p
Kai kurie net lipa ant \VUgras{vežimo} \textsuperscript{(37)}, kad sudėtų pėdus,
kuriuos jiems paduoda \VUgras{šakėmis} \textsuperscript{(38)}.

%%K t08-c1-10-1 p
10. --- Visoje \VUgras{lygumoje} \textsuperscript{(39)}, kur pakyla \VUgras{vieversiai}
\textsuperscript{(41)}, viešpatauja judrumas. Visi užsiėmę pjūtimi. Bet kol vargšai
toliau naudoja senuosius įrankius --- \VUgras{dalgius, pjautuvus} ir
\VUgras{spragilus} \textsuperscript{(40)} ---, turtingi savininkai savo kviečius ar
savo \VUgras{rugius} \textsuperscript{(43)} pjauna \VUgras{pjaunamąja mašina}
\textsuperscript{(42)}. Paskui, nereikėdami vežti pėdų į kluoną, jie krauna didelius
\VUgras{pėdų kūgius} \textsuperscript{(44)}.

%%K t08-c1-11-1 p
11. --- Tada atvaroma \VUgras{kuliamoji} \textsuperscript{(47)}, kurią
\VUgras{lokomobilis} \textsuperscript{(45)} suka \VUgras{pavaros diržu}
\textsuperscript{(46)}, ir taip grūdas atskiriamas nuo \VUgras{šiaudų,} nepaliekant
lauko, kuriame jis buvo pasėtas.

%%K t08-tit-4 sub
\VUsaut{5.0mm}
\VUpk{10.2pt}{40}{Audra.}
\VUsaut{3.4mm}

%%K t08-c1-12-1 p
12. --- Dažnai per pjūties metą kyla \VUgras{audros} \textsuperscript{(53)}.
Pirmiausia iš tolo girdėti \VUgras{griaustinio} dundesiai; ima pūsti \VUgras{vakarų
vėtra} \textsuperscript{(54)} ir, alsuodama, lenkia storiausius \VUgras{miško}
\textsuperscript{(49)} medžius. Juodi \VUgras{debesys} \textsuperscript{(56)} užpuola
dangų; juos perveria akinantys \VUgras{žaibo} \textsuperscript{(55)} zigzagai. Ima
kristi \VUgras{lietus,} o kartais ir \VUgras{kruša,} guldydami derlių, jau parengtą
pjauti.

%%K t08-c1-13-1 p
13. --- Dažnai žaibas padega pastatą. Antai pernai \VUgras{vėjo malūno}
\textsuperscript{(50)} stogas buvo paverstas pelenais, o du jo \VUgras{sparnai}
\textsuperscript{(51)} sudaužyti.

%%K t08-c1-14-1 p
14. --- Po kelių valandų vėl gimsta ramybė. Ties horizontu pasirodo spindinti
\VUgras{vaivorykštė} \textsuperscript{(52)}, ir gamta atgauna savo gyvenimą,
nutrauktą keliolikai akimirkų. Kaimiečiai, kurie prisiglaudė \VUgras{trobelėse}
\textsuperscript{(48)}, vėl imasi darbo ir drąsiai stengiasi atitaisyti žalą, kurią ką
tik patyrė.

%%K t08-tit-5 sub
\VUsaut{6.0mm}
\VUcentre{13.2pt}{90}{\textit{Antroji scena.}}
\VUsaut{3.6mm}

%%K t08-tit-6 sub
\VUpk{10.2pt}{40}{Medžioklė. --- Vynuogių nuėmimas.}
\VUsaut{1.4mm}
\VUpk{10.2pt}{40}{Žvejyba.}
\VUsaut{4.0mm}

%%K t08-c2-01-1 p
1. --- Apie \VUgras{rugpjūčio} pabaigą ar \VUgras{rugsėjo} vidurį, pagal sritis,
prasideda medžioklės atidarymas. Vos pirmieji saulės spinduliai išvijo
\VUgras{driežus} \textsuperscript{(2)} iš jų urvo, kai \VUgras{medžiotojas}
\textsuperscript{(5)} leidžiasi į kelią su savo \VUgras{šunimis}
\textsuperscript{(4)}.

%%K t08-c2-02-1 p
2. --- Jis apsivilko savo tvirtą \VUgras{medžioklės apdarą} \textsuperscript{(9)} iš
storos gelumbės, apsiavė odiniais \VUgras{antblauzdžiais,} su kuriais galės eiti
drėgna žole, kur slepiasi \VUgras{rupūžės} \textsuperscript{(1)}; jis pasiėmė savo
\VUgras{šautuvą} \textsuperscript{(6)} ir savo \VUgras{medžioklės krepšį}
\textsuperscript{(10)}, ir štai jis išeina.

%%K t08-c2-03-1 p
3. --- Vaikymosi įkarštis atveda jį į vynuogyno vidurį, kur vynuogių nuėmimas labai
gyvas. Vynuogininko vaikai, kurie paprastai gano ožkas prie kelio, šiandien palieka
jas ganytis kaip pakliuvo ir, pririšę prie kuolo seną \VUgras{ožį} \textsuperscript{(3)},
kuris nepraleistų progos pasinaudoti laisve ir pabėgti, patys taip pat pasiima savo
dalį malonumo. Vienas nutveria vynuogių kekę iš \VUgras{derliaus krepšio}
\textsuperscript{(12)} ir linksminasi spausdamas \VUgras{sultis} \textsuperscript{(14)}
į \VUgras{taurelę} \textsuperscript{(13)}, kad pasidarytų misos (nesurūgusio vyno). Kitas
puošiasi \VUgras{lapų vainiku} \textsuperscript{(15)}, kurį ką tik nupynė.

Šalia jų begėdžiai \VUgras{žvirbliai} \textsuperscript{(16)} ateina net prieš
spaustuvą pagrobti vynuogių.

%%K t08-c2-04-1 p
4. --- Kol vaikai linksminasi, vyrai dirba. \VUgras{Vynuogių skynėjai}
\textsuperscript{(18)} neša vynuoges į rūsį \VUgras{nugaros krepšiais}
\textsuperscript{(17)}; \VUgras{kubilius} \textsuperscript{(19)} taiso
\VUgras{statines} \textsuperscript{(20)}, tikrina, ar \VUgras{šulai}
\textsuperscript{(22)} ir \VUgras{dugnai} \textsuperscript{(21)} sveiki, ir keičia
sulūžusius \VUgras{lankus} \textsuperscript{(23)}. Jis taip pat taisė
\VUgras{kubilus} \textsuperscript{(25)}, į kuriuos pilamos \VUgras{vynuogės}
\textsuperscript{(26)}, kad rūgtų.

%%K t08-c2-05-1 p
5. --- Kai rūgimas baigiasi, vynas nupilamas, o likučiai dedami ant
\VUgras{spaustuvo} \textsuperscript{(28)}, ir pamažu spaudžiant \VUgras{sraigtu}
\textsuperscript{(29)} išspaudžiamos sultys, kurios teka į \VUgras{kubiliuką}
\textsuperscript{(32)}, ir dar gaunama \VUgras{vyno} \textsuperscript{(31)}.

%%K t08-tit-8 sub
\VUsaut{6.0mm}
\VUpk{10.2pt}{40}{Alus ir sidras.}
\VUsaut{3.4mm}

%%K t08-c2-06-1 p
6. --- Ant \VUgras{rūsio} \textsuperscript{(27)} sienos kopia \VUgras{apynių}
\textsuperscript{(30)} šaka. Ten jis tėra puošmena, bet kai kuriuose kraštuose, kur
vynmedis labai retas, apynys auginamas \VUgras{alui} daryti.

%%K t08-c2-07-1 p
7. --- Mažoji \VUgras{obelis} \textsuperscript{(33)} apkrauta obuoliais, kurie,
subrendę, duos puikaus \VUgras{sidro.} Savo \VUgras{narvelyje}
\textsuperscript{(34)} \VUgras{kanarėlė} \textsuperscript{(35)} ir \VUgras{dagilis}
\textsuperscript{(36)} pavydžiomis akimis žiūri į žvirblius, kurie laisvai
išdykauja.

%%K t08-c2-08-1 p
8. --- Kiek akys užmato, ant kalvų šlaito išrikiuoti \VUgras{vynmedžių kuolai}
\textsuperscript{(40)}, kurie laiko puikius \VUgras{kamienus} \textsuperscript{(39)},
nusėtus sunkiomis \VUgras{kekėmis.}

%%K t08-c2-08-2 p
Visus metus \VUgras{vynuogininkas} \textsuperscript{(42)} dirbo prižiūrėdamas savo
\VUgras{vynuogyną} \textsuperscript{(38)}, ir dabar už savo triūsą atlyginamas gražiu
derliumi. \VUgras{Vynuogių skynėjos} \textsuperscript{(41)} pripildo kubilus,
sustatytus ant vežimo, traukiamo \VUgras{mulų} \textsuperscript{(43)}, ir visi
džiaugsmingi.

%%K t08-tit-9 sub
\VUsaut{5.0mm}
\VUpk{10.2pt}{40}{Žvejyba.}
\VUsaut{3.4mm}

%%K t08-c2-09-1 p
9. --- Taip pat ir \VUgras{meškeriotojai} \textsuperscript{(44)}, kurie nuo ryto
atėjo įsitaisyti prie vandens kranto su savo \VUgras{meškerėmis}
\textsuperscript{(45)}.

%%K t08-c2-09-2 p
\VUgras{Žuvis} \textsuperscript{(47)} kimba į \VUgras{kabliuką,} vos tik jie
panardina savo \VUgras{valą} \textsuperscript{(46)} į vandenį, ir vos jie spėja
atnaujinti \VUgras{masalą,} kai nauja auka blaškosi valo gale.

%%K t08-c2-09-3 p
Vis dėlto \VUgras{tinklinis žvejys} \textsuperscript{(48)}, kuris iš savo
\VUgras{valties} \textsuperscript{(49)} meta \VUgras{užmetamąjį tinklą} į
\VUgras{upės} \textsuperscript{(50)} vidurį, pagauna daugiau žuvų.

%%K t08-c2-10-1 p
10. --- Ruduo yra gerų pasivaikščiojimų metas: pėsčiomis, \VUgras{raitomis}
\textsuperscript{(52)}, \VUgras{dviračiu} \textsuperscript{(53)},
\VUgras{automobiliu} \textsuperscript{(54)}. \VUgras{Keliai} \textsuperscript{(51)}
švarūs, ir naudojamasi paskutinėmis gražiomis dienomis, nes štai jau
\VUgras{kregždės} \textsuperscript{(56)} renkasi ant stogų, kad visais sparnais
išskristų į šiltesnius kraštus. Senojo \VUgras{riešutmedžio} \textsuperscript{(57)}
lapija geltonuoja, \VUgras{riešutai} krinta, o aukšti \VUgras{medžiai,} sustatyti
būriu kaip \VUgras{puokštė} \textsuperscript{(58)} ant \VUgras{aukštumos}
\textsuperscript{(59)}, netrukus nusimes lapus.

%%K t08-c2-11-1 p
11. --- \VUgras{Saulė} \textsuperscript{(60)} teka vėliau, dienos trumpėja,
\VUgras{rytą} ima justi gana aštrus šaltukas, ir štai jau \VUgras{slankų}
\textsuperscript{(61)} pulkai palieka mūsų nesvetingus kraštus.

\VUsaut{6.0mm}
\VUfilet{22mm}
